\chapter{Background}
\label{chap:background}

\section{Data Exploration}
\textit{Data Exploration} can be defined as the process through which users
formulate information needs over a dataset even when they do not know exactly
what they were looking for, involving the act of investingating among data,
often realizing that something is interesting only after finding it. Typically,
an exploration session includes several queries, where the result of each one
triggers the formulation of the next one, uncovering unkown or unanticipated
insights in the data, letting information that is needed to emerge and evolve
itself rather than being fully known in advance.  

As stated in \cite{idreos2015overview}, traditional data management systems
assume that users, when posing a query, have a good knowledge of the data schema
and that the query they provide leads to a result that is actually what they
wanted, presupposing that the end user already knows what they are looking for.
These assumptions are strong: not only they require the user to have a precise
understanding of the schema and semantics of the underlying data, but they
also presuppose that the user has the technical expertise needed to express
even complex queries. Moreover, distributed and heterogeneous data sources
aggravate this situation even further. On one hand, they increase the cognitive
load required of the user, who needs to know not only which data is available,
but also where it resides, a difficulty that becomes even more problematic in
data exploration processes, where we may potentially not know in advance what
information we are looking for. On the other hand, they increase the technical
competencies required to access the data itself, as different sources may expose
different query languages and technologies, such as SQL and NoSQL, each
requiring its own specific expertise simply to access the data at all.

In organizational contexts, however, restricting data exploration just to
personnel used to those query technologies can be a limitation which may exclude
domain experts who, despite lacking a technical background, may hold knowledge
that is highly relevant to bringing value to the work being carried out. For
these users, the requirement to master query technologies becomes an obstacle
not only to data exploration, but more broadly to any attempt to consult
data that is not made available to them or not in the form they would need by
the services provided to them. This situation is commonly addressed in practice
through mediation: whenever a custom, ad-hoc query is required, who needs those
data must rely on the company or department responsible for the software product
to implement it on their behalf. 

This solution, however, requires a certain amount of effort, especially in
organizations where data is distributed across different services owned by
different teams, a common situation in companies that adopt microservice
architecture, adding a coordination effort on top. This process becomes
increasingly costly, not only for the labor required from technical personnel,
but also because it imposes a considerable maintenance burden on development teams,
repeatedly involved in implementing one-off queries that fall outside the
scope of the system's predefined interfaces, diverting development time away
from work that would otherwise be more valuable to the system.
An even more critical consequence, however, concerns timing: this mediation
process introduces significant latency between the formulation of an
information need and its fulfillment, a delay that is particularly
degrading to data exploration. A further factor central to data exploration
is, in fact, interactive time: the delay between the formulation of an
information need and the availability of an answer, or at least a partial
one, sufficient to provide the user with enough context to formulate their
next request. This responsiveness is crucial, as the data exploration process
must remain interactive enough to keep the user engaged and avoid breaking
their train of thought. Distribution, in this regard, introduces additinal
latency, arising not only from the need for different nodes to communicate with
one another, but also from the fact that additional work is required to
coordinate the execution across the distributed environment and to reconcile the
partial results into a single answer. This complexity is unavoidable if we want
to enable a data exploration process that remains as transparent as possible
with respect to how the system is actually structured, a distinction that is
fundamental to support data exploration beyond the problem of data fragmentation
in today's distributed systems, particularly within microservices architectures.

Given the issues discussed in this section, this thesis aims to address
them through the design of \textit{DEX}, guided by a set of concrete objectives.\\
First, \textit{DEX} seeks to abstract users from the schema of the underlying data
sources and from the knowledge of which specific data source a given piece of data
originates from, information which is not relevant to someone engaged in data
exploration and is primarily interested in the data itself rather than in
how or where it is stored. What matters to the user is the data and its meaning,
not whether it is stored in a relational, document-oriented, key-value, graph,
or any other type of data store. \textit{DEX}'s second objective is to prevent
users from the requirement of being proficient in the specific query technology
adopted by each data source, regardless of whether it exposes a structured query
language, such as \textit{SQL} or a \textit{NoSQL} query language, or a
different access paradigm, such as the APIs typically exposed by key-value
stores. Third, \textit{DEX} aims to relieve users from the effort required to
aggregate data retrieved from potentially different sources, providing them
instead with a single, unified response: if this aggregation effort were
left to the user performing the data exploration, it would become a
significant burden, hindering an effective exploration process. \\
Given the centrality of interactive time in data exploration to enable a fluid
user experience, \textit{DEX} further strives to return a valuable response as
quickly as possible, even if the response is temporarily incomplete, providing
users with feedback that can inform and facilitate their exploration process
toward their final objective. \\ 
Taken together, these aims amount to abstracting the user from the complexity
introduced in distributed systems, while also removing when possible the need
for mediation with a technical team, freeing them to focus on work of greater
value to the system and avoiding the intrinsic latency of that mediation
process. When such mediation cannot be entirely avoided, these aims are
nonetheless pursued through a design that keeps the framework easily
maintainable and extensible, able to accommodate new data sources, new
technologies, and changes to the data made available for exploration.
This design is further guided by the goal of keeping the cognitive load imposed
on the user low, allowing data to be discovered, understood, and related to
other data with ease. \\
The overall goal of \textit{DEX} is therefore to facilitate data exploration,
all that while trying to remain efficient in terms of both system load and query
performance, making organization's data more broadly accessible to those who can 
leverage it to bring value.
